\section{Appendix Y: Boundary Conditions and Well-Posedness (Audit V5.2)}
\label{app:boundary}

To satisfy the V5 "Cauchy Problem" requirement, we explicitly define the boundary conditions used in solving the TRXT field equations.

\subsection{Y.1: Spatial Boundary Conditions}
\begin{enumerate}
    \item \textbf{Vacuum Asymptotics ($r \to \infty$):} The field must approach the uniform vacuum condensate value:
    \begin{equation}
        \lim_{r \to \infty} \Phi(t, \mathbf{x}) = \rho_0 e^{i \theta_\infty}
    \end{equation}
    where $\rho_0$ is the ground state density determined by the gap equation, and $\theta_\infty$ is a constant global phase.
    \item \textbf{Core Regularity ($r \to 0$):} To avoid singularities at the origin of spherically symmetric systems:
    \begin{equation}
        \frac{\partial \Phi}{\partial r} \bigg|_{r=0} = 0
    \end{equation}
\end{enumerate}

\subsection{Y.2: Well-Posedness Analysis}
The equation of motion $\Box_g \Phi = 0$ is a wave equation derived from the acoustic metric $g_{\mu\nu}$. The system is hyperbolic (stable evolution) if and only if the metric signature is Lorentzian $(-+++)$.
\begin{itemize}
    \item The time component $g_{00} = -1/c_s^2$.
    \item Since stability analysis (Appendix K) proves $c_s^2 > 0$, the signature is preserved.
    \item Thus, the Cauchy problem is \textbf{well-posed} for all initial data within the stability bound.
\end{itemize}
